%%%%%%%%%%%%%%%%%%%%%%%%%%%%%%%%%%%%%%%%%%%%%%%%%%%%%%%%%%%%%%%%%%%%%%%%%%%%
% CAPÍTULO 9: METODOLOGÍA
%%%%%%%%%%%%%%%%%%%%%%%%%%%%%%%%%%%%%%%%%%%%%%%%%%%%%%%%%%%%%%%%%%%%%%%%%%%%

\chapter{Metodología}
\label{ch:metodologia}

\section{Enfoque metodológico}

El TFM se aborda desde un enfoque de \textit{Design Science Research} (DSR), adecuado para trabajos de ingeniería informática en los que se diseña, construye y evalúa un artefacto software para resolver un problema concreto. Según Hevner \cite{Hevner2007}, la DSR se organiza en torno a tres ciclos complementarios (Tabla~\ref{tab:dsr}).

\begin{table}[H]
  \centering
  \begin{tabularx}{\textwidth}{Y{0.5793} Y{1.2000} Y{1.2207}}
    \hline
    \textbf{Ciclo} & \textbf{Descripción} & \textbf{Aplicación al TFM} \\
    \hline
    Ciclo de relevancia & Conectar el entorno de aplicación (problema, requisitos y criterios de aceptación) con el diseño, y devolver a ese entorno el artefacto para su prueba & Movilidad universitaria de la UAL (actuación n.º 43 del Plan Director Smart Campus UAL \cite{UALSmartCampus2019}); requisitos, casos de uso y criterios de éxito \\
    Ciclo de diseño & Iterar entre la construcción del artefacto y su evaluación & Arquitectura, modelo digital, \textit{matching} y simulación; validación funcional, técnica y experimental \\
    Ciclo de rigor & Fundamentar el diseño en la base de conocimiento (teorías, métodos y trabajos previos) y aportar a ella los resultados & Revisión sistemática (77 trabajos), ISO/IEC 25010 y métodos de análisis de simulación; contribuciones del Capítulo~\ref{ch:contribucion} \\
    \hline
  \end{tabularx}
  \caption{Ciclos de la DSR.}
  \label{tab:dsr}
\end{table}

Este enfoque permite justificar que el TFM no es solo una implementación, sino una propuesta que se evalúa y se contrasta con la literatura.

\section{Fases metodológicas}

Las fases metodológicas siguen el proceso de seis actividades de la DSR propuesto por Peffers et al. \cite{Peffers2007} y se corresponden con la revisión sistemática previa y con las fases F1--F9 del plan de trabajo (Capítulo~\ref{ch:cronograma}), tal como recoge la Tabla~\ref{tab:fases-metodologicas}.

\begin{table}[H]
  \centering
  \begin{tabularx}{\textwidth}{Y{0.8615} Y{0.4308} Y{1.6923} Y{1.0154}}
    \hline
    \textbf{Fase DSR} & \textbf{Fase Plan} & \textbf{Actividad principal} & \textbf{Resultado} \\
    \hline
    1. Identificación del problema & Previa (SLR) & Revisión sistemática, definición de RQ & Estado del arte + RQ \\
    2. Definición de objetivos & F1 & Objetivos, alcance, casos de uso, requisitos & Requisitos y diseño preliminar \\
    3. Diseño y desarrollo & F1--F6 & Arquitectura (F1); \textit{backend}, \textit{frontend}, \textit{matching}, DT y simulación (F2--F6) & Artefacto software \\
    4. Demostración & F6--F7 & Ejecución de escenarios y casos de prueba & Resultados de simulación y pruebas \\
    5. Evaluación & F7 & Comparación de escenarios, análisis de métricas & Resultados y discusión \\
    6. Comunicación & F8--F9 & Redacción de la memoria y preparación de la defensa & Memoria final + presentación \\
    \hline
  \end{tabularx}
  \caption{Fases metodológicas (DSR \texorpdfstring{$\leftrightarrow$}{<->} plan).}
  \label{tab:fases-metodologicas}
\end{table}

\section{Técnicas de validación}

La validación se articula en tres niveles complementarios (Tabla~\ref{tab:validacion}).

\begin{table}[H]
  \centering
  \begin{tabularx}{\textwidth}{Y{0.3704} Y{1.6296}}
    \hline
    \textbf{Nivel} & \textbf{Detalles} \\
    \hline
    1 --- Validación funcional & Técnica: pruebas unitarias, de integración y de sistema. Objetivo: comprobar que los casos de uso (CU-01 a CU-13) funcionan según lo especificado. Herramientas: Jest para el \textit{backend}, Vitest para el \textit{frontend}, Cypress para las pruebas de extremo a extremo y Postman para las APIs. \\
    2 --- Validación técnica & Técnica: pruebas de rendimiento y verificación de sincronización. Objetivo: comprobar que el \textit{matching} responde en tiempos aceptables y que el Gemelo Digital refleja los cambios de estado. Métricas: tiempo de respuesta del \textit{matching}, latencia de sincronización, errores por escenario. \\
    3 --- Validación experimental & Técnica: simulación de escenarios controlados (S0--S5). Objetivo: analizar el efecto del carpooling sobre indicadores de movilidad. Métricas: vehículos utilizados, vehículos evitados, ocupación media, solicitudes cubiertas, distancia adicional. \\
    \hline
  \end{tabularx}
  \caption{Técnicas de validación.}
  \label{tab:validacion}
\end{table}

\section{Rigor y trazabilidad}

Para garantizar el rigor metodológico:
\begin{itemize}
    \item Cada objetivo específico (OE1--OE12) se traza contra al menos una RQ y una fase de validación.
    \item Cada requisito funcional (RF) se traza contra un caso de uso, y cada requisito no funcional (RNF) contra una característica del modelo ISO/IEC 25010.
    \item Cada escenario de simulación se traza contra un conjunto de métricas.
    \item Los resultados se documentan con datos reproducibles (semillas, configuraciones, versiones).
\end{itemize}

\section{Limitaciones metodológicas}

\begin{itemize}
    \item La validación experimental se basa en datos sintéticos, no en datos reales de movilidad de la UAL.
    \item La evaluación con usuarios reales queda fuera del alcance del TFM.
    \item Los resultados de simulación representan el comportamiento del modelo bajo condiciones definidas, no predicciones del mundo real.
\end{itemize}
